\documentclass{amsart}
\usepackage{amssymb}

\newtheorem{theorem}{Theorem}[section]
\newtheorem{lemma}[theorem]{Lemma}
\newtheorem{proposition}[theorem]{Proposition}
\newtheorem{corollary}[theorem]{Corollary}
\newtheorem{conjecture}[theorem]{Conjecture}
\theoremstyle{remark}
\newtheorem{remark}[theorem]{Remark}

\numberwithin{equation}{section}

\DeclareMathOperator{\des}{des}
\newcommand{\RR}{\mathbb{R}}
\newcommand{\CC}{\mathbb{C}}
\newcommand{\HH}{\mathbb{H}}
\newcommand{\fS}{\mathfrak{S}}
\renewcommand{\Im}{\operatorname{Im}}
\newcommand{\cU}{\mathcal{U}}
\newcommand{\cT}{\mathcal{T}}

\hyphenation{manuscript}

\begin{document}

\title[Real zeros of $A_n(x)+kxA_{n-2}(x)$]{Real zeros of $A_n(x)+kxA_{n-2}(x)$ and a conjecture of Zhang}

\author{Knowledge Sense}
\address{Knowledge Sense}
\email{math@knowledgesense.jp}

\date{9 October 2026}

\subjclass[2020]{Primary 05A15, 26C10; Secondary 05A05, 30C15}

\keywords{Eulerian polynomial, real-rootedness, interlacing, tangent number, stack sortable permutation}

\begin{abstract}
Let $A_n(x)$ be the Eulerian polynomial of the symmetric group $\fS_n$. We
determine, for every $n\ge 3$, the set of real numbers $k$ for which
$A_n(x)+kxA_{n-2}(x)$ has $n-1$ distinct real zeros: this happens exactly
when $k<-n(n-1)$ or $k>-T_{m+1}/T_m$, where $m=\lfloor n/2\rfloor$ and the
$T_m$ are the tangent numbers. This confirms a conjecture of Zhang. We also
locate all zeros for every real $k$. At $k=-n(n-1)$ the number $1$ is a
double zero, and at $k=-T_{m+1}/T_m$ the number $-1$ is a multiple zero. For
$k$ strictly between these values exactly two zeros are nonreal, and they are
complex conjugates of modulus one. The proof expresses $A_n$ through
polynomials $F_n$ obtained from a constant by alternately applying two fixed
first-order differential operators, and shows by elementary arguments that
the zeros of $F_{n-2}$ and $F_n$ strictly interlace.
\end{abstract}

\maketitle

\section{Introduction}

For $n\ge 1$ let
\[
A_n(x)=\sum_{\pi\in\fS_n}x^{\des(\pi)},
\]
where $\des(\pi)$ is the number of descents of $\pi$, and put $A_0(x)=1$.
Thus $A_1=1$, $A_2=1+x$, $A_3=1+4x+x^2$, and $\deg A_n=n-1$ for $n\ge 1$. For
$n\ge 3$ and $k\in\RR$ we consider
\begin{equation}\label{eq:K}
K_{n,k}(x)=A_n(x)+kxA_{n-2}(x).
\end{equation}

Polynomials of this form arise from stack sorting. B\'ona \cite{Bona} conjectured
that the descent polynomials of $t$-stack sortable permutations have only
real zeros, and for $t=n-2$ the descent polynomial is $K_{n,-1}(x)$; see
\cite[Section 1]{Zhang}. Br\"and\'en \cite[Theorem 5.2]{Branden} proved that $K_{n,k}$ has only real and simple zeros
for all $n\ge 3$ and $k>-2$; in \cite{Branden} the Eulerian polynomials carry an extra
factor $x$, which does not affect the statement. Zhang \cite{Zhang} gave another
proof based on the $\mathbf{s}$-Eulerian polynomials of Savage and Visontai
\cite{SavageVisontai}, and showed that $K_{n,k}$ has only real zeros for $n>3$ and $k\ge -n$
\cite[Theorem 1.3]{Zhang}; the restriction $n>3$ is needed there, since
$K_{3,-3}(x)=1+x+x^2$. Stanley suggested determining when $K_{n,k}$ has only
real zeros \cite[Section 3]{Zhang}. Zhang first considered the case of
distinct real zeros, determined the range of $k$ numerically for
$3\le n\le 18$, and proposed the following conjecture.

Let the tangent numbers $T_m$ be defined by
\begin{equation}\label{eq:tangent}
\tan t=\sum_{m\ge 1}T_m\frac{t^{2m-1}}{(2m-1)!},
\end{equation}
so that $T_1=1$, $T_2=2$, $T_3=16$, $T_4=272$, and $T_m$ is the number of
alternating permutations of $\{1,\dots,2m-1\}$. This is the normalization
used in \cite{Zhang}.

\begin{conjecture}[{\cite[Conjecture 3.2]{Zhang}}]\label{conj:Zhang}
For any $n\ge 3$, the polynomial $K_{n,k}(x)$ has only real and distinct
zeros if and only if $k\in(-\infty,-n(n-1))\cup(-a(\lfloor n/2\rfloor),
+\infty)$, where $a(m)=T_{m+1}/T_m$.
\end{conjecture}

Since $K_{n,k}$ has degree $n-1$ for every $k$, having only real and
distinct zeros means having $n-1$ distinct real zeros. Our main result
confirms Conjecture~\ref{conj:Zhang}.

\begin{theorem}\label{thm:main}
Let $n\ge 3$, $m=\lfloor n/2\rfloor$ and $k\in\RR$. Then $A_n(x)+kxA_{n-2}(x)$
has $n-1$ distinct real zeros if and only if
\[
k<-n(n-1)\qquad\text{or}\qquad k>-\frac{T_{m+1}}{T_m}.
\]
\end{theorem}

Theorem~\ref{thm:main} is contained in the following description of all
zeros.

\begin{theorem}\label{thm:zeros}
Let $n\ge 3$, $m=\lfloor n/2\rfloor$ and $\tau_n=T_{m+1}/T_m$. Then
$\tau_n<n(n-1)$. For every $k\in\RR$ the polynomial $K_{n,k}$ has degree
$n-1$ and satisfies $K_{n,k}(0)=1$, and the following hold.
\begin{enumerate}
\item[(a)] If $k>-\tau_n$, then all zeros of $K_{n,k}$ are negative and
simple.
\item[(b)] If $k=-\tau_n$, then $-1$ is a zero of multiplicity two if $n$ is
odd and of multiplicity three if $n$ is even; all other zeros are negative
and simple.
\item[(c)] If $-n(n-1)<k<-\tau_n$, then $K_{n,k}$ has exactly two nonreal
zeros; they are simple, complex conjugate and of modulus one. All other
zeros are negative and simple.
\item[(d)] If $k=-n(n-1)$, then $1$ is a zero of multiplicity two; all other
zeros are negative and simple.
\item[(e)] If $k<-n(n-1)$, then $K_{n,k}$ has exactly two positive zeros,
both simple; all other zeros are negative and simple.
\end{enumerate}
In particular, $K_{n,k}$ has only real zeros if and only if
$k\le -n(n-1)$ or $k\ge -\tau_n$.
\end{theorem}

The last assertion of Theorem~\ref{thm:zeros} answers the question of
Stanley mentioned above. For $3\le n\le 18$ the numbers $-n(n-1)$ and $-\tau_n$ are the values
$\omega_n$ and $\Omega_n$ tabulated in \cite[Section 3]{Zhang}; for instance
$\tau_8=\tau_9=T_5/T_4=7936/272=496/17$. Zhang also observed that the
products $\prod_{i=1}^{j}\Omega_{2i+1}$ are tangent numbers up to sign. This
is explained by Theorem~\ref{thm:zeros}, since
$\prod_{i=1}^{j}\tau_{2i+1}=\prod_{i=1}^{j}T_{i+1}/T_i=T_{j+1}$.

\subsection*{Method}
Since $A_n$ is palindromic, it can be written as
\[
A_n(x)=(x+1)^{e_n}x^{d_n}H_n(x+x^{-1})
\]
with a monic polynomial $H_n$ in the variable $w=x+x^{-1}$
(Lemma~\ref{lem:Hn}). It is convenient to work
with the variable $z=x/(x-1)^2$, which satisfies $z^{-1}=w-2$. In this
variable the recurrence of the Eulerian polynomials becomes the alternating
application of the two operators
\[
\cU f=f+zf',\qquad \cT f=(1+6z)f+z(1+4z)f'
\]
to a constant (Lemma~\ref{lem:Fn}). Neither operator depends on $n$, so
$F_n$ and $F_{n-2}$ are mapped to $F_{n+1}$ and $F_{n-1}$ by the same
operator. A logarithmic-derivative argument shows that $\cU$ and $\cT$
preserve real-rootedness, and the Hermite--Kakeya--Obreschkoff theorem, in
the special case proved in Section~\ref{sec:interlace}, then propagates
strict interlacing from $(F_{n-2},F_n)$ to $(F_{n-1},F_{n+1})$
(Proposition~\ref{prop:interlace}). Consequently the ratio $H_n/H_{n-2}$ is
increasing to the right of the largest zero of $H_{n-2}$, and only one
quadratic factor of $K_{n,k}$ can leave the negative real axis. The two
thresholds are the negatives of the values of this ratio at $w=2$ and
$w=-2$, that is, at $x=1$ and $x=-1$. They are computed from $A_n(1)=n!$
and the exponential generating function of the Eulerian polynomials.
Br\"and\'en's proof of \cite[Theorem 5.2]{Branden} likewise changes
variables, by the substitution $x\mapsto x/(1+x)$, and uses a two-step
recursion.

\subsection*{Related results}
Combinations of two consecutive Eulerian polynomials have been studied from
the point of view of Hurwitz stability. Yang and Zhang
\cite[Theorem 2.3]{YangZhang} showed that $(x+1)A_n(x)+kxA_{n-1}(x)$ has no
zeros in the open right half-plane for $n\ge 2$ and $k\ge -n$; in
\cite{YangZhang} the Eulerian polynomials are indexed differently. They also
conjectured, in terms of the Euler zigzag numbers, the exact range of $k$
for which this polynomial has no zeros in the closed right half-plane, and
asked for a connection with Conjecture~\ref{conj:Zhang}
\cite[Conjecture 3.2 and the remark following it]{YangZhang}. Ding and Zhu
\cite[Proposition 2.17]{DingZhu} gave a criterion of this kind for
polynomials satisfying a general recurrence. With $\beta=n$,
$\gamma=\phi=1$, $\nu=-1$, $\rho=1$ and $\eta=k$ in their notation, it shows
that $(x+1)A_{n+1}(x)+kxA_n(x)$ is Hurwitz stable, in their sense of having
no zeros in the open right half-plane, for $n\ge 0$ and $k\ge -(n+1)$. The
present paper treats the pencil~\eqref{eq:K}, which involves $A_n$ and
$A_{n-2}$, and does not address the question of Yang and Zhang.

\subsection*{Organization}
Section~\ref{sec:F} introduces the polynomials $F_n$ and $H_n$. Section
\ref{sec:interlace} proves the strict interlacing of $F_{n-2}$ and $F_n$.
Section~\ref{sec:proof} evaluates the thresholds and proves
Theorems~\ref{thm:main} and~\ref{thm:zeros}.

\section{The polynomials $F_n$ and $H_n$}\label{sec:F}

We use the classical recurrence
\begin{equation}\label{eq:rec}
A_{n+1}(x)=(1+nx)A_n(x)+x(1-x)A_n'(x)\qquad(n\ge 0)
\end{equation}
and the exponential generating function
\begin{equation}\label{eq:egf}
\sum_{n\ge 0}A_n(x)\frac{t^n}{n!}=\frac{x-1}{x-e^{(x-1)t}};
\end{equation}
see \cite{Comtet}. For completeness we note that, for the series
$G=\sum_{n\ge 0}A_n(x)t^n/n!$, the recurrence \eqref{eq:rec} is equivalent to
\[
(1-xt)G_t=G+x(1-x)G_x,
\]
and together with $G(x,0)=1$ this equation determines the coefficients of
$G$ recursively. A direct computation shows that the right-hand side of
\eqref{eq:egf} satisfies the same equation and equals $1$ at $t=0$.

For $n\ge 1$ put
\[
d_n=\Big\lfloor\frac{n-1}{2}\Big\rfloor,\qquad e_n=n-1-2d_n\in\{0,1\}.
\]
Thus $e_n=0$ for odd $n$ and $e_n=1$ for even $n$. For $f\in\RR[z]$ define
\[
\cU f=f+zf',\qquad \cT f=(1+6z)f+z(1+4z)f',
\]
and define polynomials $F_n\in\RR[z]$ by $F_1=1$ and
\[
F_{2m}=\cU F_{2m-1}\quad(m\ge 1),\qquad F_{2m+1}=\cT F_{2m}\quad(m\ge 1).
\]
For example $F_2=1$, $F_3=1+6z$, $F_4=1+12z$ and $F_5=1+30z+120z^2$.

\begin{lemma}\label{lem:Fn}
For every $n\ge 1$,
\[
A_n(x)=(x+1)^{e_n}(x-1)^{2d_n}F_n\!\left(\frac{x}{(x-1)^2}\right).
\]
\end{lemma}

\begin{proof}
The identity is between rational functions of $x$. Put
$\varphi(x)=x/(x-1)^2$. Then
\begin{equation}\label{eq:phi}
\varphi'=-\frac{x+1}{(x-1)^3},\qquad
1+4\varphi=\frac{(x+1)^2}{(x-1)^2},\qquad
1+6\varphi=\frac{x^2+4x+1}{(x-1)^2}.
\end{equation}
We use induction on $n$; for $n=1$ both sides equal $1$. Assume the identity
for $n$, and write $F$ and $F'$ for $F_n(\varphi)$ and $F_n'(\varphi)$.

Let $n=2m+1$ with $m\ge 0$, so that $A_n=(x-1)^{2m}F$. By \eqref{eq:phi},
\[
x(1-x)A_n'=-2mx(x-1)^{2m}F+x(x+1)(x-1)^{2m-2}F',
\]
and \eqref{eq:rec} gives
\begin{align*}
A_{n+1}&=(x+1)(x-1)^{2m}F+x(x+1)(x-1)^{2m-2}F'\\
&=(x+1)(x-1)^{2m}\bigl(F+\varphi F'\bigr).
\end{align*}
This is the claim for $n+1=2m+2$, since $F_{2m+2}=\cU F_{2m+1}$,
$e_{2m+2}=1$ and $d_{2m+2}=m$.

Let $n=2m$ with $m\ge 1$, so that $A_n=(x+1)(x-1)^{2m-2}F$. By
\eqref{eq:phi},
\begin{align*}
x(1-x)A_n'&=-x(x-1)^{2m-1}F-(2m-2)x(x+1)(x-1)^{2m-2}F\\
&\quad+x(x+1)^2(x-1)^{2m-4}F'.
\end{align*}
Adding $(1+2mx)A_n$, the coefficient of $(x-1)^{2m-2}F$ becomes
\[
(1+2mx)(x+1)-(2m-2)x(x+1)-x(x-1)=x^2+4x+1,
\]
and by \eqref{eq:phi},
\[
A_{n+1}=(x-1)^{2m}\bigl((1+6\varphi)F+\varphi(1+4\varphi)F'\bigr).
\]
This is the claim for $n+1=2m+1$, since $F_{2m+1}=\cT F_{2m}$,
$e_{2m+1}=0$ and $d_{2m+1}=m$.
\end{proof}

A real polynomial is \emph{standard} if its leading coefficient is positive.
We write $I$ for the open interval $(-1/4,0)$.

\begin{lemma}\label{lem:roots}
For every $n\ge 1$, $F_n$ is standard of degree $d_n$, $F_n(0)=1$, and $F_n$
has $d_n$ simple zeros, all in $I$.
\end{lemma}

\begin{proof}
Both operators fix the constant term, so $F_n(0)=1$. If $f$ has degree $d$
and leading coefficient $a$, then $\cU f$ has degree $d$ and leading coefficient
$(d+1)a$, and $\cT f$ has degree $d+1$ and leading coefficient $(4d+6)a$. This
gives the degree and the sign of the leading coefficient by induction.

It remains to show that $\cU$ and $\cT$ map a standard polynomial $f$ of degree
$d$ with $d$ simple zeros in $I$ to a polynomial whose zeros are simple and
lie in $I$. If $d=0$, then $\cU f$ is a positive constant and $\cT f=f\cdot(1+6z)$
has the single zero $-1/6\in I$. Let $d\ge 1$ and let $r_1<\dots<r_d$ be the
zeros of $f$. Since $f$ is standard, $f'(r_j)$ has the sign $(-1)^{d-j}$,
$f(0)>0$, and $f(-1/4)$ has the sign $(-1)^d$. Since
$r_j<0$ and $1+4r_j>0$, both
\[
\cU f(r_j)=r_jf'(r_j)\qquad\text{and}\qquad \cT f(r_j)=r_j(1+4r_j)f'(r_j)
\]
have the sign $(-1)^{d-j+1}$, while $\cU f(0)=\cT f(0)=f(0)>0$. Hence each of $\cU f$
and $\cT f$ has a zero in each of the $d$ intervals $(r_j,r_{j+1})$ with
$1\le j<d$ and $(r_d,0)$. This accounts for all zeros of $\cU f$. Moreover
$\cT f(-1/4)=-f(-1/4)/2$ has the sign $(-1)^{d+1}$, whereas $\cT f(r_1)$ has the
sign $(-1)^d$, so $\cT f$ has one more zero in $(-1/4,r_1)$. This accounts for
all $d+1$ zeros of $\cT f$, and all of them lie in $I$.
\end{proof}

For $n\ge 1$ define
\[
H_n(w)=(w-2)^{d_n}F_n\!\left(\frac{1}{w-2}\right).
\]

\begin{lemma}\label{lem:Hn}
For every $n\ge 1$, $H_n$ is a monic polynomial of degree $d_n$,
\begin{equation}\label{eq:Hn}
A_n(x)=(x+1)^{e_n}x^{d_n}H_n(x+x^{-1}),
\end{equation}
and the zeros of $H_n$ are the numbers $2+\rho^{-1}$, where $\rho$ runs over
the zeros of $F_n$. In particular, $H_n$ has $d_n$ simple zeros, all less
than $-2$.
\end{lemma}

\begin{proof}
Write $d=d_n$ and $F_n(z)=\sum_{i=0}^{d}c_iz^i$. Then
$H_n(w)=\sum_{i=0}^{d}c_i(w-2)^{d-i}$, which is monic of degree $d$ because
$c_0=F_n(0)=1$. Since $1/\varphi(x)=x+x^{-1}-2$, we have
$(x-1)^{2d}F_n(\varphi(x))=x^d\varphi(x)^{-d}F_n(\varphi(x))
=x^dH_n(x+x^{-1})$, and \eqref{eq:Hn} follows from Lemma~\ref{lem:Fn}. As
$H_n(2)=c_d\ne 0$, a number $w$ is a zero of $H_n$ if and only if $w\ne 2$
and $(w-2)^{-1}$ is a zero of $F_n$. The map $\rho\mapsto 2+\rho^{-1}$ is
injective and sends $I$ into $(-\infty,-2)$, so the last statement follows
from Lemma~\ref{lem:roots}.
\end{proof}

\section{Strict interlacing}\label{sec:interlace}

Let $f,g\in\RR[z]$ with $\deg f=d\ge 1$ and $\deg g=d-1$. We say that $g$
\emph{strictly interlaces} $f$, and write $g\prec f$, if $f$ has $d$ simple
real zeros $r_1<\dots<r_d$, $g$ has $d-1$ simple real zeros
$s_1<\dots<s_{d-1}$, and
\[
r_1<s_1<r_2<s_2<\dots<s_{d-1}<r_d.
\]
For $d=1$ the condition only says that $g$ is a nonzero constant. The
following two lemmas are special cases of the theorem of Hermite, Kakeya
and Obreschkoff \cite{Obreschkoff}, \cite{Dedieu}, \cite[Theorem 2.4]{Branden}; we include
proofs for completeness.

\begin{lemma}\label{lem:pencil}
Let $f$ and $g$ be standard with $g\prec f$, $\deg f=d\ge 1$, and let
$s_1<\dots<s_{d-1}$ be the zeros of $g$. Then the following hold.
\begin{enumerate}
\item[(i)] There are $c>0$, $b\in\RR$ and $c_1,\dots,c_{d-1}>0$ with
\[
\frac{f(z)}{g(z)}=cz+b-\sum_{j=1}^{d-1}\frac{c_j}{z-s_j}.
\]
\item[(ii)] The Wronskian $W=f'g-fg'$ is positive on $\RR$.
\item[(iii)] For every $t\in\RR$, the polynomial $f+tg$ has $d$ simple real
zeros. In the $d$ open intervals into which $s_1,\dots,s_{d-1}$ divide
$\RR$, it has exactly one zero each, and $f/g$ is strictly increasing on each
of them.
\end{enumerate}
\end{lemma}

\begin{proof}
(i) Since $\deg f=\deg g+1$ and the zeros of $g$ are simple, partial
fractions give such an expansion with $c$ the ratio of the leading
coefficients and with residue $-c_j=f(s_j)/g'(s_j)$ at $s_j$. The sign of
$g'(s_j)$ is $(-1)^{d-1-j}$. Exactly $d-j$ zeros of $f$ exceed $s_j$, so the
sign of $f(s_j)$ is $(-1)^{d-j}$. Hence $c_j>0$.

(ii) Away from the zeros of $g$ we have
\[
W=g^2\Bigl(\frac fg\Bigr)'=g^2\Bigl(c+\sum_j\frac{c_j}{(z-s_j)^2}\Bigr)>0,
\]
and at $s_j$ we have $W(s_j)=-f(s_j)g'(s_j)=c_jg'(s_j)^2>0$.

(iii) By (i), $f/g$ is strictly increasing on each of the $d$ intervals; on
each of them it tends to $-\infty$ at the left end and to $+\infty$ at the
right end. Hence $f/g=-t$ has exactly one solution in each interval. These
are zeros of $f+tg$, and $f+tg$ does not vanish at $s_j$ because
$f(s_j)\ne 0$. Since $\deg(f+tg)=d$, these are all its zeros.
\end{proof}

\begin{lemma}\label{lem:converse}
Let $p$ and $q$ be standard, $\deg p=d\ge 1$ and $\deg q=d-1$. Assume that
the zeros of $q$ are real and simple, that $p$ and $q$ have no common zero,
and that $p+tq$ has only real zeros for every $t\in\RR$. Then $q\prec p$.
\end{lemma}

\begin{proof}
If $d=1$ there is nothing to prove, so let $d\ge 2$ and let
$s_1<\dots<s_{d-1}$ be the zeros of $q$. Let $\HH$ be the open upper
half-plane. The function $\psi=p/q$ is holomorphic on $\HH$. It takes no real
value there: if $\psi(z)=-t$ with $t\in\RR$ and $z\in\HH$, then $z$ would be a
nonreal zero of $p+tq$. Since $\HH$ is connected, $\Im\psi$ has constant sign
on $\HH$. As $y\to+\infty$ we have $\psi(iy)=c\,iy+O(1)$, where $c>0$ is the
ratio of the leading coefficients, so $\Im\psi>0$ on $\HH$.

Partial fractions give $\psi(z)=cz+b+\sum_j h_j/(z-s_j)$ with $b,h_j\in\RR$,
and $h_j=p(s_j)/q'(s_j)\ne 0$ because $p$ and $q$ have no common zero. As
$\varepsilon\to 0^+$,
\[
\Im\psi(s_j+i\varepsilon)=-\frac{h_j}{\varepsilon}+O(1),
\]
so $h_j<0$. Hence $\psi$ has the form in Lemma~\ref{lem:pencil}(i). The
argument in the proof of Lemma~\ref{lem:pencil}(iii), with $t=0$, shows that $p$ has exactly one zero in
each of the intervals $(-\infty,s_1)$, $(s_j,s_{j+1})$ and
$(s_{d-1},\infty)$. These are $d$ simple real zeros, and $q\prec p$.
\end{proof}

\begin{lemma}\label{lem:preserve}
If $f\in\RR[z]$ is nonzero and has only real zeros, then $\cU f$ and $\cT f$
are nonzero and have only real zeros.
\end{lemma}

\begin{proof}
Let $d=\deg f$ and let $a$ be the leading coefficient of $f$. The leading
coefficients of $\cU f$ and $\cT f$ are $(d+1)a$ and $(4d+6)a$, as in the
proof of Lemma~\ref{lem:roots}, so both are nonzero. Write
$f=a\prod_j(z-r_j)$ with real $r_j$, repeated according to multiplicity.
Let $z\notin\RR$. Then $z$, $1+4z$ and $f(z)$ are nonzero, and
since $(1+6z)/(z(1+4z))=z^{-1}+\frac12(z+\frac14)^{-1}$,
\begin{align*}
\frac{\cU f(z)}{zf(z)}&=\frac1z+\sum_j\frac{1}{z-r_j},\\
\frac{\cT f(z)}{z(1+4z)f(z)}&=\frac1z+\frac{1/2}{z+1/4}+\sum_j\frac{1}{z-r_j}.
\end{align*}
For real $r$ we have $\Im (z-r)^{-1}=-\Im z/|z-r|^2$, so the imaginary part
of each right-hand side is $-\Im z$ times a positive number. Hence
$\cU f(z)\ne 0$ and $\cT f(z)\ne 0$.
\end{proof}

\begin{proposition}\label{prop:interlace}
For every $n\ge 3$ we have $F_{n-2}\prec F_n$ and $H_{n-2}\prec H_n$, and all
zeros of $H_n$ and $H_{n-2}$ are less than $-2$.
\end{proposition}

\begin{proof}
We prove $F_{n-2}\prec F_n$ by induction on $n$. For $n=3$ it holds because
$F_1=1$ and $F_3=1+6z$. Let $n\ge 3$ and assume $F_{n-2}\prec F_n$. Put
$L=\cU$ if $n$ is odd and $L=\cT$ if $n$ is even. By the definition of the $F_j$,
\[
F_{n+1}=LF_n\qquad\text{and}\qquad F_{n-1}=LF_{n-2}.
\]
We write $L=u+v\,\frac{d}{dz}$; thus
\[
(u,v)=(1,z)\ \text{if}\ L=\cU,\qquad (u,v)=\bigl(1+6z,\,z(1+4z)\bigr)\ \text{if}\ L=\cT.
\]

Both $F_n$ and $F_{n-2}$ are standard by Lemma~\ref{lem:roots}. By
Lemma~\ref{lem:pencil}(iii), $F_n+tF_{n-2}$ is nonzero and has only real
zeros for every $t\in\RR$, so by Lemma~\ref{lem:preserve} the same is true of
$L(F_n+tF_{n-2})=F_{n+1}+tF_{n-1}$. By Lemma~\ref{lem:roots}, $F_{n+1}$ and
$F_{n-1}$ are standard, $\deg F_{n+1}=d_{n+1}=d_{n-1}+1=\deg F_{n-1}+1$, and
the zeros of $F_{n-1}$ are simple and lie in $I$.

Suppose that $\zeta$ is a common zero of $F_{n+1}$ and $F_{n-1}$. Then
$\zeta\in I$, so $v(\zeta)\ne 0$. For any $f$ and $g$,
\[
(Lf)g-(Lg)f=v\,(f'g-fg').
\]
Applied to $f=F_n$ and $g=F_{n-2}$ at $z=\zeta$, this gives
$v(\zeta)W(\zeta)=0$, where $W=F_n'F_{n-2}-F_nF_{n-2}'$. This contradicts
Lemma~\ref{lem:pencil}(ii). Hence $F_{n+1}$ and $F_{n-1}$ have no common zero,
and Lemma~\ref{lem:converse} gives $F_{n-1}\prec F_{n+1}$.

Now let $n\ge 3$ and $d=d_n$. By Lemma~\ref{lem:Hn}, the zeros of $H_n$ and
$H_{n-2}$ are simple and less than $-2$, and they are the images of the
zeros of $F_n$ and $F_{n-2}$ under the map $\rho\mapsto 2+\rho^{-1}$, which
is strictly decreasing on $I$. If $d=1$, then $H_{n-2}=1$ and $H_n$ has a
single simple real zero, so $H_{n-2}\prec H_n$. If $d\ge 2$, let
$\rho_1<\dots<\rho_d$ be the zeros of $F_n$ and $\sigma_1<\dots<\sigma_{d-1}$
those of $F_{n-2}$. Since $F_{n-2}\prec F_n$, their images satisfy
\[
2+\rho_d^{-1}<2+\sigma_{d-1}^{-1}<2+\rho_{d-1}^{-1}<\dots<2+\sigma_1^{-1}
<2+\rho_1^{-1}<-2.
\]
Thus $H_{n-2}\prec H_n$.
\end{proof}

\section{Proof of the main theorems}\label{sec:proof}

\begin{lemma}\label{lem:values}
We have $H_n(2)=n!/2^{e_n}$ for $n\ge 1$. Moreover,
$H_{2m+1}(-2)=T_{m+1}$ for $m\ge 0$, and $H_{2m}(-2)=T_{m+1}/2$ for
$m\ge 1$.
\end{lemma}

\begin{proof}
At $x=1$ we have $x+x^{-1}=2$, and $A_n(1)=n!$, so \eqref{eq:Hn} gives the
first claim.

At $x=-1$, \eqref{eq:egf} becomes
\[
\sum_{n\ge 0}A_n(-1)\frac{t^n}{n!}=\frac{2}{1+e^{-2t}}=1+\tanh t.
\]
Since $\tanh t=-i\tan(it)$, \eqref{eq:tangent} gives
$\tanh t=\sum_{m\ge 0}(-1)^mT_{m+1}t^{2m+1}/(2m+1)!$. Hence
$A_{2m}(-1)=0$ for $m\ge 1$ and $A_{2m+1}(-1)=(-1)^mT_{m+1}$ for $m\ge 0$.

For $n=2m+1$, \eqref{eq:Hn} reads $A_{2m+1}(x)=x^mH_{2m+1}(x+x^{-1})$, and at
$x=-1$ it gives $(-1)^mT_{m+1}=(-1)^mH_{2m+1}(-2)$.

For $n=2m$ with $m\ge 1$, \eqref{eq:Hn} reads
$A_{2m}(x)=(x+1)x^{m-1}H_{2m}(x+x^{-1})$. Differentiating at $x=-1$, where the
factor $x+1$ vanishes, gives
\[
A_{2m}'(-1)=(-1)^{m-1}H_{2m}(-2).
\]
On the other hand, \eqref{eq:rec} at $x=-1$ with $n=2m$ gives
\[
A_{2m+1}(-1)=(1-2m)A_{2m}(-1)-2A_{2m}'(-1)=-2A_{2m}'(-1).
\]
Therefore
$(-1)^mT_{m+1}=2(-1)^mH_{2m}(-2)$.
\end{proof}

\begin{proof}[Proof of Theorem~\ref{thm:zeros}]
Fix $n\ge 3$ and write $d=d_n\ge 1$, $e=e_n$ and $m=\lfloor n/2\rfloor$.
Since $d_{n-2}=d-1$ and $e_{n-2}=e$, \eqref{eq:Hn} gives
$xA_{n-2}(x)=(x+1)^ex^dH_{n-2}(x+x^{-1})$, and therefore
\begin{equation}\label{eq:KQ}
K_{n,k}(x)=(x+1)^ex^dQ_k(x+x^{-1}),\qquad Q_k=H_n+kH_{n-2}.
\end{equation}
The leading coefficient of $A_n$ is $1$ and $\deg xA_{n-2}=n-2$, so
$K_{n,k}$ has degree $n-1$; also $K_{n,k}(0)=A_n(0)=1$.

Let $R=H_n/H_{n-2}$. Since $H_{n-2}\prec H_n$ by
Proposition~\ref{prop:interlace}, Lemma~\ref{lem:pencil}(iii) with $t=k$
shows that the polynomial $Q_k$ is monic of degree $d$ and has $d$ simple real zeros
$w_1<\dots<w_d$. If $d=1$, then $H_{n-2}=1$, $R=H_n$ is a monic linear
polynomial, and we put $J=\RR$; the only zero of $Q_k$ is the solution $w_1$
of $R(w)=-k$. If $d\ge 2$, let $s_{d-1}$ be the largest zero of $H_{n-2}$ and
put $J=(s_{d-1},\infty)$. By Lemma~\ref{lem:pencil}(iii), the zeros
$w_1,\dots,w_{d-1}$ are less than $s_{d-1}$, and $w_d$ is the unique solution
of $R(w)=-k$ in $J$; moreover $s_{d-1}<-2$ by
Proposition~\ref{prop:interlace}, so $w_1,\dots,w_{d-1}$ are less than $-2$.

In both cases $R$ is continuous and strictly increasing on $J$ by
Lemma~\ref{lem:pencil}(i), the interval $J$ contains $-2$ and $2$, and $w_d$
is the unique point of $J$ with $R(w_d)=-k$. Consequently, for
$a\in\{-2,2\}$,
\begin{equation}\label{eq:compare}
\begin{aligned}
w_d<a&\iff k>-R(a),\\
w_d=a&\iff k=-R(a),\\
w_d>a&\iff k<-R(a).
\end{aligned}
\end{equation}
By Lemma~\ref{lem:values}, and because $e_{n-2}=e_n$,
\[
R(2)=\frac{n!}{(n-2)!}=n(n-1).
\]
If $n=2m+1$, then $R(-2)=H_{2m+1}(-2)/H_{2m-1}(-2)=T_{m+1}/T_m$; if $n=2m$,
then $R(-2)=(T_{m+1}/2)/(T_m/2)=T_{m+1}/T_m$. Thus $R(-2)=\tau_n$ for both
parities, and $\tau_n=R(-2)<R(2)=n(n-1)$ because $R$ is increasing on $J$.

Since $x^dQ_k(x+x^{-1})=\prod_{j=1}^{d}x\,(x+x^{-1}-w_j)$, equation
\eqref{eq:KQ} becomes
\begin{equation}\label{eq:factor}
K_{n,k}(x)=(x+1)^e\prod_{j=1}^{d}q_{w_j}(x),\qquad q_w(x)=x^2-wx+1.
\end{equation}
The zeros of $q_w$ have product $1$ and sum $w$. Hence $q_w$ has two distinct
negative zeros if $w<-2$, equals $(x+1)^2$ if $w=-2$, has two complex
conjugate nonreal zeros of modulus one if $-2<w<2$, equals $(x-1)^2$ if
$w=2$, and has two distinct positive zeros if $w>2$. A zero $x$ of $q_w$
determines $w=x+x^{-1}$, so $q_w$ and $q_{w'}$ have no common zero when
$w\ne w'$, and $-1$ is a zero of $q_w$ only for $w=-2$.

Since $w_1,\dots,w_{d-1}$ are distinct and less than $-2$, the factors
$q_{w_1},\dots,q_{w_{d-1}}$ contribute $2(d-1)$ distinct negative zeros,
none equal to $-1$. If $e=1$, the factor $x+1$ contributes the zero $-1$.
These zeros are simple zeros of $(x+1)^e\prod_{j<d}q_{w_j}$, and none of them
is a zero of $q_{w_d}$ unless $w_d=-2$, in which case $q_{w_d}=(x+1)^2$.
Therefore all zeros of $K_{n,k}$ outside those of $q_{w_d}$ are negative and
simple, and \eqref{eq:compare} with $R(-2)=\tau_n$ and $R(2)=n(n-1)$ gives
the five cases.

If $k>-\tau_n$, then $w_d<-2$ and the zeros of $q_{w_d}$ are negative,
simple and different from $-1$; this is (a). If $k=-\tau_n$, then $w_d=-2$,
$q_{w_d}=(x+1)^2$, and $-1$ has multiplicity $2+e$; this is (b). If
$-n(n-1)<k<-\tau_n$, then $-2<w_d<2$; this is (c). If $k=-n(n-1)$, then
$w_d=2$ and $q_{w_d}=(x-1)^2$; this is (d). If $k<-n(n-1)$, then $w_d>2$ and
$q_{w_d}$ has two distinct positive zeros; this is (e). The final statement
follows from (a) to (e).
\end{proof}

\begin{proof}[Proof of Theorem~\ref{thm:main}]
By Theorem~\ref{thm:zeros}, $K_{n,k}$ has $n-1$ distinct real zeros in cases
(a) and (e), and it has a multiple or a nonreal zero in cases (b), (c) and
(d).
\end{proof}

\begin{remark}
The proof also shows how the zeros move with $k$. For $d\ge 2$ the numbers
$w_1,\dots,w_{d-1}$ always lie below the largest zero of $H_{n-2}$, so only
the two zeros of $q_{w_d}$ can leave the negative real axis. As $k$
decreases from $+\infty$ to $-\infty$, the parameter $w_d$ increases from
$s_{d-1}$ (from $-\infty$ if $d=1$) to $+\infty$. The two zeros of $q_{w_d}$
move along the negative axis until they meet at $-1$, where $K_{n,k}$ then
has a zero of multiplicity $2+e_n$. They then move along the unit circle
until they meet at $1$, where $K_{n,k}$ has a double zero, and finally they
separate along the positive axis.
\end{remark}

\section*{Use of AI-assisted tools}

AI-assisted tools were used in the research and preparation of this manuscript.

\begin{thebibliography}{9}

\bibitem{Bona}
M. B\'ona, \emph{Symmetry and unimodality in $t$-stack sortable
permutations}, J. Combin. Theory Ser. A \textbf{98} (2002), no.~1, 201--209.

\bibitem{Branden}
P. Br\"and\'en, \emph{On linear transformations preserving the P\'olya
frequency property}, Trans. Amer. Math. Soc. \textbf{358} (2006), no.~8,
3697--3716. Cited from arXiv:math/0403364v2; theorem numbers refer to
this version.

\bibitem{Comtet}
L. Comtet, \emph{Advanced combinatorics}, D. Reidel Publishing Co.,
Dordrecht, 1974.

\bibitem{Dedieu}
J.-P. Dedieu, \emph{Obreschkoff's theorem revisited: what convex sets are
contained in the set of hyperbolic polynomials?}, J. Pure Appl. Algebra
\textbf{81} (1992), no.~3, 269--278.

\bibitem{DingZhu}
M.-J. Ding and B.-X. Zhu, \emph{Stability of combinatorial polynomials and
its applications}, preprint (2021), arXiv:2106.12176v2.

\bibitem{Obreschkoff}
N. Obreschkoff, \emph{Verteilung und Berechnung der Nullstellen reeller
Polynome}, VEB Deutscher Verlag der Wissenschaften, Berlin, 1963.

\bibitem{SavageVisontai}
C. D. Savage and M. Visontai, \emph{The $\mathbf{s}$-Eulerian polynomials
have only real roots}, Trans. Amer. Math. Soc. \textbf{367} (2015), no.~2,
1441--1466.

\bibitem{YangZhang}
A. L. B. Yang and P. B. Zhang, \emph{The real-rootedness of Eulerian
polynomials via the Hermite--Biehler theorem}, preprint (2015),
arXiv:1501.05824v2.

\bibitem{Zhang}
P. B. Zhang, \emph{On the real-rootedness of the descent polynomials of
$(n-2)$-stack sortable permutations}, Electron. J. Combin. \textbf{22}
(2015), no.~4, Paper 4.12.

\end{thebibliography}

\end{document}
